\chapter{Research Scope and Requirements}
\label{chap:motivation}

Chapter~\ref{chap:evaluating-agents} illustrates why evaluating the coordinator's routing decision is hard, particularly when the goal is to catch a wrong decision before the coordinator delegates the request to a sub-agent. Section~\ref{sec:cost-of-attribution} shows that attributing a failure to a specific agent and step by inspecting a completed trace is costly and does not scale when performed by hand. Section~\ref{sec:failure-attribution} shows that automating this attribution with an LLM-as-a-judge is not much more reliable: on the Who\&When benchmark, the best automated method reached 53.5\% accuracy at identifying the responsible agent and 14.2\% accuracy at pinpointing the responsible step, with some methods scoring below random \cite{zhangWhichAgentCauses2025}.

A check that runs at the moment of the coordinator's decision, before the selected sub-agent is invoked, would attribute the failure to the coordinator itself and could prevent the rest of the trace from running. A source for that check could come from probing the hidden states of the coordinator's underlying LLM.

This chapter reviews the literature on hidden-state probing for routing decisions and identifies the gaps it leaves open. The chapter then determines the requirements that guide the rest of this work.

\section{Designs for a Routing Check}
\label{sec:routing-check-designs}

To attribute a wrong decision to the coordinator, and to prevent the wrong routing from taking effect, a check needs to be available before the coordinator invokes the selected sub-agent.

One option is a deterministic check, which applies a heuristic or a policy, encoded as a rule, over the request. Section~\ref{subsec:building-approaches} makes the same point about a hand-authored specification in general, in the context of classical agents. Writing it correctly demands domain expertise, and it is incomplete by construction, since it covers only the cases it was written to anticipate. A rule matched against a request is this kind of specification, so its coverage is bounded the same way, and it does not generalize across a sub-agent catalog it was not written for.

A language model can also serve as this check, reused as a judge that grades the coordinator's routing decision. As Section~\ref{sec:llm-as-judge} shows, a judge model carries systematic biases of its own. It also adds latency and cost, since it requires a separate model call for every decision it checks. Moreover, because the judge is itself an LLM performing a similar reasoning task, it is subject to the same class of failure as the coordinator it is meant to check.

A classifier trained to predict the correct sub-agent label from a fixed representation of the request's text can be a further source for evaluation. The representation is computed without the routing labels, and only the classifier that reads it is trained on them. This work refers to a classifier of this kind as a \emph{text-only classifier}. It can be trained on lexical features, counting how many times each word occurs in a query (a \emph{bag-of-words} representation), optionally weighted by how rare the word is across the training set using \emph{TF-IDF}. An alternative is a classifier trained on \emph{sentence embeddings}: dense vectors produced by an embedding model that capture the meaning of a sentence. In this case, two queries phrased differently but with the same meaning can map to nearby vectors. A text-only classifier of any of these three kinds can be more efficient than an LLM, and can also exceed a zero-shot LLM's accuracy on the same routing task \cite{balaMultiAgentRoutingSetValued2026}. However, this is not guaranteed: depending on the routing task, an LLM may perform better than a text-only classifier. Where a text-only classifier is more accurate than the coordinator's own LLM, its disagreement with the coordinator's decision could flag a likely misrouting, or the classifier could replace the coordinator's classification entirely, rather than only check it.

Wu, Zhu et al. showed that a probe, defined in Section~\ref{subsec:probing-method}, can capture a decision of this kind directly from a model's hidden states, with high accuracy: the tool a model is about to call is linearly encoded in its hidden states, and readable before the call is generated (Section~\ref{subsec:readability}) \cite{wuToolCallingLinearly2026}. A routing decision selects one sub-agent from a declared set, the same kind of selection that Wu, Zhu et al. studied for tool calls, so the information a check needs may already be present in the coordinator's hidden states before it generates a response. A probe trained on these hidden states could therefore be a stronger source than a text-only classifier for this decision, classifying the correct sub-agent directly from them and serving as an evaluation flag before the coordinator's decision takes place. However, this probe is beneficial only if its accuracy exceeds that of a text-only classifier: otherwise, it offers no advantage over a text-only classifier, which costs less to train and does not require access to hidden states.


\section{Related Work on Hidden-State Probing for Routing}
\label{sec:related-work-routing}

\subsection{Hidden State Probing for Routing and Evaluation}
\label{subsec:hidden-state-probing}

Hidden-state probing has been applied to classification tasks. One line of this work has focused on exploiting hidden states to choose which general-purpose LLM, from a fixed pool, should answer a query. Results report that a probe trained on the hidden states can outperform a classifier trained on the query's surface text.

Varshney et al. trained a probe on prefill activations, the hidden states from the forward pass over the prompt before generation starts, to predict, for each candidate model, its probability of answering a query correctly. The probe outperformed baseline classifiers built from sentence embeddings, retrieval, and text-complexity features, reaching a mean \emph{AUROC} (area under the ROC curve, a threshold-independent measure of how well a classifier ranks positive cases above negative ones) of 0.856 against 0.804 for the strongest baseline on the frontier model pool \cite{varshneyLLMRouterRethinking2026}.

Lugoloobi et al. reported the same pattern for a model's own success on mathematics and coding tasks: a probe trained on activations read before generation outperformed a TF-IDF baseline, with the margin ranging from 0.05 to 0.20 AUROC across models and tasks \cite{lugoloobiLLMsEncodeTheir2026}.

A separate study by Wu, Wang et al. found that a router trained on cross-layer hidden states outperformed routers built on output probabilities, embeddings, and verbose self-evaluation, for choosing between a local and a cloud model, with relative improvements of 16.68\% in AUROC and 18.86\% on the high-accuracy deployment scenario over the best baseline \cite{wuFairComprehensiveEvaluation2026}.

Reading a decision from the hidden states before generation has also been applied to tasks beyond choosing which model should answer.

Chen et al. built two classifiers on hidden states read during prompt prefill, to classify a prompt's topic as math, code, or general text: a trained linear probe, and a training-free classifier that reduces the hidden state to summary statistics and compares them to a reference statistic computed for each candidate topic. Direct LLM inference on this same task was less reliable than the trained probe, particularly for smaller models under zero-shot prompting: its accuracy improved only with larger models and in-context learning, needed careful prompt design to reach a reasonable level, and failed outright for the smallest model tested. The training-free classifier trailed the LLMs' own decisions on several cases, but generalized better to out-of-distribution and ambiguous prompts \cite{chenFastIntentClassification2025}.

Cencerrado et al. trained a linear probe on activations read after a question but before generation, to predict whether a model's own forthcoming answer would be correct. On out-of-distribution questions, the probe generalized better than an embedding-based baseline and a baseline built from \emph{verbalized confidence} (a confidence score the model states in its own text output), though the embedding-based baseline matched or slightly exceeded the probe on in-distribution questions \cite{cencerradoNoAnswerNeeded2026}.

Healy et al. trained a classifier on an agent's final-layer hidden states, taken from the same forward pass that generated a tool call. The classifier flagged whether that call hallucinated the wrong function, malformed its parameters, or bypassed the available tools, reaching up to 86.4\% detection accuracy \cite{healyInternalRepresentationsIndicators2026}.

Sun et al. trained a linear probe on an LLM's hidden state at the last input token, read before generation, to predict whether a task needs a tool call at all. The probe reached an AUROC of 0.89 to 0.96 across six models, exceeding the same models' own verbalized reasoning about tool necessity \cite{sunLLMAgentsAlready2026}.

Across the works reviewed in this subsection, hidden states carry information about a decision beyond what the model's own output reveals. Because this information is available before generation, a probe reading it may supply a signal that can be compared against the model's eventual output. This motivates using a probe as an evaluation mechanism that flags a routing error before it takes effect.

\subsection{LatentGate}
\label{subsec:latentgate}

Ratnakar et al. came closest to routing among specialized destinations with a hidden-state probe. LatentGate is a router that reads hidden states from a frozen small language model and classifies a query directly against a pool of specialized agents, to reduce the latency and cost of a prompt-based LLM router \cite{ratnakarLatentGateLowLatencySemantic2026}. The paper reports that a prompt-based LLM router takes 1,500 to 2,000 milliseconds per query, with cost scaling with the number of agents, against about 28 milliseconds for LatentGate, and estimates a 99.98\% reduction in the monthly cost of routing 1 million queries \cite{ratnakarLatentGateLowLatencySemantic2026}. LatentGate reaches this lower cost because it stops the forward pass at a fixed interior layer of the small language model and never generates output tokens. A linear probe then classifies the pooled hidden state at that layer into one of the candidate agents. The frozen model never receives the list of candidate agents, and it never receives a routing instruction. The query alone is tokenized and passed through the model, under a maximum sequence length of 128 tokens \cite{ratnakarLatentGateLowLatencySemantic2026}, so that its hidden states serve as a representation of the query for the probe to classify.

LatentGate is evaluated on its own synthetic benchmark of 100 enterprise agents, built so that agents in the same domain share vocabulary on purpose. On this benchmark's out-of-distribution queries, paraphrased and deliberately built to overlap in vocabulary with other agents, LatentGate reaches 80.0\% accuracy, against 60.0\% for a sentence-embedding classifier, BGE-M3 \cite{ratnakarLatentGateLowLatencySemantic2026}. Because the model is never instructed to route, its own routing accuracy is never measured, so the probe is never compared against the model it reads from. LatentGate attributes about 65\% of its own errors on its benchmark to the deliberate cross-domain vocabulary overlap \cite{ratnakarLatentGateLowLatencySemantic2026}, but the paper never reports the accuracy of a bag-of-words classifier on this benchmark, so whether such a classifier fails for the same reason, and by how much, is not established. The only such baseline that LatentGate reports is on CLINC150, a separate 150-class intent-classification benchmark. There, LatentGate outperforms an SVM and TF-IDF classifier by 1.7 percentage points, 92.9\% against 91.2\% \cite{ratnakarLatentGateLowLatencySemantic2026}. The small accuracy difference suggests that most of the routing signal on that benchmark is already recoverable from surface text, so the probe's contribution is limited. LatentGate's scope is also limited to single-turn queries \cite{ratnakarLatentGateLowLatencySemantic2026}.

\subsection{What the Literature Leaves Open}
\label{subsec:literature-gap}

The hidden-state probing work reviewed in Section~\ref{subsec:hidden-state-probing} and Section~\ref{subsec:latentgate} evaluates every probe on single-turn queries. None of it tests whether a hidden-state probe still routes correctly once a query follows more than one turn of conversation. Furthermore, none of it compares a probe's routing accuracy against the model's own routing accuracy, the comparison that defines a probe's effectiveness as an evaluation mechanism.

As Section~\ref{subsec:latentgate} shows, LatentGate is the reviewed work closest in approach, but its purpose differs, because it produces the routing decision instead of checking a decision the coordinator has already made. It is also the only one that builds its benchmark to be hard for a sentence-embedding classifier, by giving agents in the same domain overlapping vocabulary on purpose, and then attributes its own errors to that overlap. It does not run the same check for a bag-of-words classifier. LatentGate's scale also differs from a coordinator's: it routes among 100 agents, not the roughly three to seven a production coordinator typically selects among, as Section~\ref{sec:multi-agent} shows. At that scale, a classifier can fail because there are many candidates to choose among, not because any two are hard to tell apart. LatentGate's benchmark does not establish whether the same difficulty holds when the candidate pool is small and the only source of difficulty is overlap between a few agents.


\section{From the Literature Gap to the Research Questions}
\label{sec:rq-derivation}

Section~\ref{subsec:literature-gap} shows that no reviewed work evaluates a probe as an evaluation mechanism checked against the model's own decision, across both single-turn and multi-turn settings. Each research question below draws on one or more of these gaps.

\subsection{Probing as an Evaluation Mechanism}
\label{subsec:rq1-derivation}

\phantomsection\label{rq:rq1}
\begin{rqbox}[RQ1]
On requests where the correct sub-agent is difficult to determine from the wording alone, does a probe on the hidden states of a coordinator model recover the correct sub-agent more accurately than the coordinator's own selection?
\end{rqbox}


\noindent The probes reviewed in Section~\ref{sec:related-work-routing} either replace the routing decision outright, as LatentGate does, or check a model's own success on a task unrelated to routing, as Healy et al. and Cencerrado et al. do. None evaluates whether a probe can serve as a check on a coordinator's own routing choice, compared against the coordinator's own accuracy. The comparison targets requests where the wording does not determine the correct sub-agent. When the wording determines it, the probe's accuracy is explained by the request text, and the comparison against the coordinator does not isolate the contribution of the hidden states.

\subsection{Comparison with the Text-Only Classifiers}
\label{subsec:rq2-derivation}

\phantomsection\label{rq:rq2}
\begin{rqbox}[RQ2]
Does a probe on the coordinator's hidden states recover the correct sub-agent more accurately than classifiers trained on a fixed representation of the query text, such as raw word counts, TF-IDF weighting, or sentence embeddings?
\end{rqbox}


\noindent \hyperref[rq:rq1]{\textbf{RQ1}} compares the probe against the coordinator. This question compares the probe against the text-only classifiers, on the same requests. A probe can serve as an evaluation instrument only if it outperforms the text-only classifiers by a clear margin, on a dataset confirmed to be hard for those classifiers. Otherwise, a text-only classifier is sufficient.

\subsection{Consistency of the Optimal Layer}
\label{subsec:rq3-derivation}

\phantomsection\label{rq:rq3}
\begin{rqbox}[RQ3]
Does the optimal probe layer fall at a consistent relative depth across different LLMs?
\end{rqbox}


\noindent Comparing a probe against the coordinator and against the text-only classifiers requires training the probe at every layer of every model tested, since no single layer is known in advance to carry the routing decision best. This setup raises a further question: whether the best-performing layer falls at the same position across different models, or varies. The models tested differ in layer count, so the position is expressed as a fraction of network depth to make it comparable. If that relative depth is consistent, an evaluation probe could be trained on a fixed subset of layers, without a separate search to determine the optimal layer for each new model.

\subsection{Disagreement as a Misrouting Flag}
\label{subsec:rq4-derivation}

\phantomsection\label{rq:rq4}
\begin{rqbox}[RQ4]
Can the disagreement between the probe and the coordinator be used as a runtime flag for likely misrouting?
\end{rqbox}


\noindent This question asks whether the disagreement between the probe and the coordinator can serve as a runtime flag for likely misrouting. Such a flag could signal a trajectory likely to fail before the selected sub-agent runs. It could stop the run early, or prompt the coordinator to re-evaluate its routing choice in response to the disagreement.


\section{Requirements for a New Dataset}
\label{sec:dataset-requirements}

Answering the research questions defined in Section~\ref{sec:rq-derivation} requires both a single-turn and a multi-turn dataset. In this work, a \emph{sample} is a single data point of a dataset: one request, together with its target label. In both datasets, each sample must be a request meant to be submitted to a coordinator agent, which must route it to one of its available sub-agents. Each sample must also carry a target label, which defines the sub-agent to which the request must be routed. The rest of this section describes each requirement.

\subsection{Routing Task}
\label{subsec:routing-setup}

The number of sub-agents used must satisfy the following requirements.

\begin{itemize}
    \item It must stay within the range Section~\ref{sec:multi-agent} reports for a production coordinator, roughly three to seven sub-agents.
    \item It must allow requests for different sub-agents to share vocabulary, so that the wording alone does not reveal which sub-agent must handle the request.
\end{itemize}

Each sample must still be routable to one, and only one, sub-agent, so that a prediction on that sample can be marked right or wrong, depending on whether the chosen sub-agent matches the correct one. The shared vocabulary must therefore hold at the level of the wording, while the content of the request still selects a single sub-agent.

\subsection{Difficulty for Text-Only Classifiers}
\label{subsec:lexical-separability}

To answer \hyperref[rq:rq2]{\textbf{RQ2}}, the dataset must be challenging for text-only classifiers. This gives the probe margin to score above them.

Three text-only classifiers are used to measure how well this difficulty holds: a logistic regression trained on bag-of-words counts, a logistic regression trained on TF-IDF-weighted counts, and a logistic regression trained on sentence embeddings. The dataset must aim for the lowest accuracy these classifiers can reach, while every sample still reads as natural and remains routable to exactly one sub-agent, without an ambiguity that would leave it unroutable.

The sub-agents must also have a balanced number of requests, and be balanced in writing style, punctuation, and sentence pattern, in addition to vocabulary. Otherwise, a text-only classifier could route a request by its class frequency or by a stylistic marker instead of by what the request asks.

\subsection{Multi-turn Setting}
\label{subsec:multiturn-setting}

The multi-turn dataset must be built on the same sub-agents as the single-turn dataset, with the same writing styles, the same number of samples, and the same balance of samples across sub-agents. The only difference is that each multi-turn sample distributes its request across more than one turn, as defined in Chapter~\ref{chap:llm} (Section~\ref{subsec:query-prompt}), instead of stating it in one.

Each multi-turn sample must use the same number of turns, so that the length of the conversation is controlled. This count must stay fixed because the dataset is meant to show how performance changes between the single-turn and the multi-turn setting, not how it changes with the number of turns.

Each multi-turn sample must start with a request and end with a request. The final request must have no response: it is the request to be routed. Together, the turns in a sample are required to carry enough information to route the complete conversation to exactly one sub-agent.


\section{Requirements for Model Selection and Probing}
\label{sec:model-selection-requirements}

\subsection{Model Pool}
\label{subsec:model-pool-requirements}

\hyperref[rq:rq3]{\textbf{RQ3}} asks whether the best-performing layer falls at a consistent relative depth across different models. To answer this research question, data must be gathered from more than one LLM, from different model families and with different parameter sizes, from a few billion to several tens of billions. The pool of models must also satisfy the following requirements.

\begin{itemize}
    \item The models must vary in layer depth, so that models with more or fewer hidden layers can be compared. Layer depth is not fixed by parameter size alone: two models of similar size can stack a different number of layers, so depth is treated as a separate axis.
    \item Each model must be instruction-tuned, as described in Section~\ref{subsec:instruction-tuning}. This training is required for the model to read the dataset's routing prompt and return a routing decision, since the coordinator's own accuracy on this decision is one of the references that a probe is evaluated against, as Section~\ref{subsec:rq1-derivation} states.
    \item Each model must support tool calling, so that the coordinator can route the request to a sub-agent by invoking it as a tool.
    \item Each model must generate with sampling disabled (greedy decoding, described in Section~\ref{subsec:decoding}), so that its routing decision is deterministic and the run can be replicated.
\end{itemize}

Each model's weights must also be publicly available, so it can be run locally with a library such as Hugging Face's Transformers, which allows the model's hidden states to be read directly from the forward pass at every layer. A closed-source model, such as the models that OpenAI or Anthropic serve, is reached only through a network API. This API returns generated text and does not return the internal hidden states.

\subsection{Probe Input Position}
\label{subsec:probe-position-requirements}

Section~\ref{subsec:last-token} explains why the position a probe reads within the prompt matters. Most of the reviewed literature reads the hidden state at the last token of the prompt, before generation: this is the position used by Sun et al., by Cencerrado et al., and by Varshney et al., described in Section~\ref{subsec:hidden-state-probing}. LatentGate instead pools the hidden states over every token of the query \cite{ratnakarLatentGateLowLatencySemantic2026}.

This work reads hidden states from the last token of the prompt. Most of the reviewed literature reports good accuracy at this position, and reading only the last token is more efficient than probing every token of the query.

\subsection{Anisotropy and Whitening}
\label{subsec:anisotropy-scope}

At the center of LatentGate's method is a correction for anisotropy (Section~\ref{subsec:anisotropy}), applied through PCA-whitening before its probe is trained. This correction spreads the representations apart along the directions where they concentrate. It removes the shared component that makes representations of different classes look similar, which makes the classes easier for a linear probe to separate and raises its accuracy \cite{ratnakarLatentGateLowLatencySemantic2026}.

This work does not apply a whitening correction like LatentGate does, since raising the probe's peak accuracy is not a goal of this work. The classifier used as a probe is trained directly on the hidden states.

